\chapter{Urbanisation in India: Patterns, Slums and Violence}

\section{The Three Stages of Urbanisation and Over-Urbanisation}

\begin{KeyConcept}
\begin{itemize}
\item \textbf{Kingsley Davis}: urbanisation is a change in \textbf{economic activities} (from agricultural to non-agricultural) plus a \textbf{demographic explosion} --- and it moves in a \textbf{cycle} with three stages.
\item The \textbf{three stages}: (1) the \textbf{initial stage} --- a rural, agricultural economy with traditional social norms and a \textbf{dispersed settlement pattern} (peasant society); (2) the \textbf{acceleration stage} --- non-agricultural economic domains develop and people (and the state) invest in \textbf{transport and communication networks}; (3) the \textbf{terminal stage} --- \textbf{over 70\%} of the population is urbanised and outside the agricultural sector.
\item In an urban society, \textbf{status shifts from land} (the dominant caste's marker) to \textbf{modern criteria} --- \textbf{well-being} and \textbf{profession} (white-collar vs blue-collar); land can even lower one's status.
\item \textbf{Over-urbanisation}: urban misery and rural poverty \textbf{exist side by side} --- some centres (Delhi/NCR/Gurgaon) are over-developed while others are neglected, producing an \textbf{imbalance}.
\end{itemize}
\end{KeyConcept}

\begin{QuickRevision}
\begin{itemize}
\item Kingsley Davis: urbanisation = change in economic activity + demographic explosion; three stages (initial/dispersed $\rightarrow$ acceleration/transport $\rightarrow$ terminal/70\% non-agricultural).
\item Status shifts from land to well-being/profession; over-urbanisation = urban misery + rural poverty side by side.
\end{itemize}
\end{QuickRevision}

\section{Amitabh Kundu and Pseudo-Urbanisation}

\begin{KeyConcept}
\begin{itemize}
\item \textbf{Amitabh Kundu} (human geography) argues India shows \textbf{pseudo-urbanisation}: people come to cities \textbf{not because the cities attract them (pull) but because of rural distress (push)} --- not industrialisation.
\item The \textbf{patterns of Indian urbanisation}: (1) \textbf{lopsided urbanisation} --- only a few \textbf{mega-cities} (mostly colonial cities or princely political capitals) develop; (2) urbanisation \textbf{without a strong economic base or industrialisation} --- population grows exponentially but the economic base \textbf{shrinks} (Delhi's industries were relocated for pollution); (3) urbanisation is mainly a product of \textbf{demographic explosion and poverty-induced rural--urban migration} (rural distress = poverty, caste discrimination, religion, ethnicity); (4) it produces \textbf{slums} (centres of poverty); (5) the migration is of \textbf{poor quality} (unskilled people learning on the job) and leads to \textbf{urban decay} (infrastructure not in sync with the numbers).
\item \textbf{Urban agglomeration} = a continuous urban spread --- a town plus its adjoining \textbf{urban outgrowths} (JNU, ports, railway colonies, cantonments). \textbf{Noida = New Okhla Industrial Development Area}.
\item \textbf{Definition of urban}: a \textbf{statutory town} (municipality/corporation/cantonment) or a \textbf{census town} (population > 5,000, density > 400, and \textbf{75\% of the male working population in non-agricultural activities} --- Indian urbanisation is largely male-led).
\end{itemize}
\end{KeyConcept}

\section{Causes and Patterns of Urbanisation in India}

\begin{KeyConcept}
\begin{itemize}
\item The \textbf{causes of urbanisation in India}: (1) \textbf{industrialisation} (Noida, Jamshedpur, Bokaro --- steel plants); (2) \textbf{commercialisation} (centres of trade --- 'smokeless capitalism' like Singapore/Hong Kong; parts of Mumbai); (3) \textbf{social benefits and services} (education, healthcare, recreation --- Allahabad, Kanpur as education centres); (4) \textbf{employment opportunities} (Bangalore as the Silicon Valley of India after liberalisation); (5) \textbf{modernisation and changes in the way of living}; (6) \textbf{rural--urban transformation} (a rural area gaining amenities).
\item \textbf{Western vs Indian urbanisation}: in the West, urbanisation is \textbf{fuelled by industrialisation} (surplus agricultural labour absorbed into industry); in India, it is \textbf{not} --- it is fuelled by rural--urban migration and demographic explosion.
\end{itemize}
\end{KeyConcept}

\begin{QuickRevision}
\begin{itemize}
\item Kundu: pseudo-urbanisation (push not pull; distress not industry); lopsided (mega-cities); no economic base; slums; poor-quality migration; urban decay.
\item Urban agglomeration/outgrowth (JNU, ports); Noida = New Okhla Industrial Development Area; statutory vs census town (5000/400/75\% male non-agricultural).
\item Causes: industrialisation, commercialisation, social services, employment (Bangalore), modernisation, rural--urban transformation; West = industry-led, India = migration-led.
\end{itemize}
\end{QuickRevision}

\section{Slums: The Census and the Sociological View}

\begin{KeyConcept}
\begin{itemize}
\item \textbf{Census 2011} defines slums as residential areas whose dwellings are \textbf{unfit for habitation} --- dilapidated, overcrowded, faulty in design, narrow streets, lacking ventilation/light/sanitation. The main differentiating factors are \textbf{sanitation, toilets and drinking water within the premises}: drinking water within the premise (56.7\% slum vs 71.2\% non-slum); \textbf{no toilet within the premises (34\% slum vs 18.6\% non-slum)}; open drainage (36.9\% slum vs 44.5\% non-slum). On most other parameters (permanent houses, home ownership, electricity, mobile phones) slums and non-slums are similar.
\item \textbf{Sociologists do not see slums as homogeneous}. \textbf{Charles Stokes} divides them into \textbf{slums of hope} and \textbf{slums of despair}:
\item \textbf{Slums of hope} --- those where people see a chance of \textbf{social mobility}, pooling resources through \textbf{kinship networks} (e.g. \textbf{Dharavi}, a manufacturing hub with expanding population and success stories). \textbf{UNFPA} found slums can be sites of mobility (e.g. Dharavi youth wrongly arrested by the police, who then studied law and became lawyers defending their area --- the story behind the film \emph{Shahid}).
\item \textbf{Slums of despair} --- \textbf{exclusionary spaces} with negligible chances of mobility, involved in deviant/illegal trades (drugs, prostitution) --- e.g. Sanjay Camp in Delhi (originally partition camps, later 'unauthorised' slums).
\end{itemize}
\end{KeyConcept}

\begin{QuickRevision}
\begin{itemize}
\item Census 2011: slums = unfit for habitation; differentiating = toilet (34\% vs 18.6\%), drinking water (56.7\% vs 71.2\%).
\item Sociologists: slums not homogeneous; Charles Stokes: slums of hope (Dharavi, mobility, kinship pooling) vs slums of despair (Sanjay Camp, deviance).
\end{itemize}
\end{QuickRevision}

\section{The Significance of Slums}

\begin{KeyConcept}
\begin{itemize}
\item The \textbf{significance of slums in Indian urbanisation}:
\item 1. \textbf{Source of cheap labour} --- Indian urbanisation is cheap (unlike the West, where labour is costly --- see the Devyani Khobragade case); slums keep wages low and are the 'engine' of urbanisation.
\item 2. \textbf{They provide urban culture} --- the 'ocean of urban culture': migrants bring folk culture (Chhath, Rath Yatra), create a \textbf{fusion culture}, and produce \textbf{counter-culture} (Mumbai's rap music; the film \emph{ABCD}) --- migration follows family $\rightarrow$ kinship $\rightarrow$ caste $\rightarrow$ village.
\item 3. \textbf{Informal business/industry} --- about \textbf{20,000 mini-factories in Dharavi}.
\item 4. \textbf{Centre of social cohesion} --- people help each other in disasters, and rebuild communities after violence (e.g. 1984 anti-Sikh violence --- studied by Veena Das).
\item 5. \textbf{Vote bank} --- a large vote bank (e.g. Manoj Tiwari and the Bihari vote bank near the Yamuna).
\item 6. \textbf{A distinct political culture} --- attracted by \textbf{welfare narratives} (free electricity/water/education --- these small savings aid mobility); Ajay Gudavarthy links Congress's decline to liberalisation ending the welfare state.
\item 7. \textbf{Engines of social change} --- through trade-union and labour movements in the city, and by carrying urban attitudes (gender equality, education, anti-superstition) back to the village.
\item \textbf{Census 2011}: states investing in social sectors have fewer slums --- \textbf{Kerala has the lowest (1.5\%)}.
\end{itemize}
\end{KeyConcept}

\section{Slums and Sexual Violence: Asis Nandy and Ayona Datta}

\begin{KeyConcept}
\begin{itemize}
\item \textbf{Mike Davis} ('Planet of Slums'): Western slums sit at the \textbf{periphery} (pushed out as the city expands), but Indian cities are \textbf{'living cities'} --- slums are interspersed with residential areas, so all classes live and interact nearby (contrast \textbf{Chandigarh/Lutyens' Delhi}, the \textbf{'dead cities'} planned only for the upper class).
\item \textbf{Asis Nandy}: sexual violence in slums is \textbf{anomic rape} --- a product of \textbf{anomic urbanisation} (Durkheim's anomie): the transition to a non-agricultural economy happened without creating the new values; \textbf{kinship ties are dying, community is weakening, the pathological becomes normal, relations become superficial}, and there is a \textbf{stranger-misogyny} and an \textbf{erosion of the public from the city} (slum-dwellers are excluded from the city's aesthetic centres).
\item \textbf{Ayona Datta} criticises the assumption that providing basic amenities will reduce crime: most sexual violence happens in the \textbf{private sphere} and goes unreported. The \textbf{lack of sanitation facilities turns women's private life into part of the public gaze} (communal toilets), and the \textbf{law/state is an enabler} of these crimes --- the state does not create the facilities for privacy, even while the courts speak of the right to privacy.
\item (Also to note: \textbf{slums as a site of intrusion/violence} --- the Sanjay Gandhi--Jagmohan slum clearances and Emergency sterilisation --- and the idea of \textbf{disposable children} from high infant mortality, are flagged for the next class.)
\end{itemize}
\end{KeyConcept}

\begin{QuickRevision}
\begin{itemize}
\item Mike Davis: Western peripheral slums vs Indian 'living cities' (slums interspersed); Chandigarh/Lutyens = 'dead cities'.
\item Asis Nandy: anomic rape; anomic urbanisation; kinship dying; stranger-misogyny; erosion of public from city.
\item Ayona Datta: private-sphere violence unreported; sanitation $\rightarrow$ public gaze; law/state as enabler.
\item Slums: cheap labour, urban culture (fusion/counter-culture), 20,000 Dharavi factories, social cohesion (1984), vote bank, welfare-narrative politics, engines of change; Kerala lowest slums (1.5\%).
\end{itemize}
\end{QuickRevision}

